\section{Full seven-product series proof}\label{sec:seven-product-proof}
For positive imaginary periods $t,s,d$ with $t+7d=2s+4$, put $P=e^{2\pi it}$, $Q=e^{2\pi is}$, $D=e^{2\pi id}$ and $Z=e^{2\pi iz}$. Then $PD^7=Q^2$. In a common convergence region the fifth-smoothed logarithms are
\begin{align*}
 \log S_5(z;a,b)&=5\sum_{\substack{n\ge1\\5\nmid n}}
 \frac{Z^n-(A_0B_0/Z)^n}{n(1-A_0^n)(1-B_0^n)},\\
 \log F_{t,s}(z)&=-5\sum_{\substack{n\ge1\\5\nmid n}}
 \frac{P^nZ^n-Q^nZ^{-n}}{n(1-P^n)(1-Q^n)},
\end{align*}
where $A_0=e^{2\pi ia},B_0=e^{2\pi ib}$ in the first line. At one surviving frequency suppress the index $n$. Summing seven shifts multiplies the two coefficients by $\sum_{k=0}^{6}D^k$ and $\sum_{k=0}^{6}D^{-k}$. The exact identities are
\begin{align*}
 \frac{P\sum_{k=0}^{6}D^k}{(1-P)(1-Q)}
 &=\frac{1+Q}{(1-P)(1-D)}-\frac1{(1-Q)(1-D)},\\
 \frac{Q\sum_{k=0}^{6}D^{-k}}{(1-P)(1-Q)}
 &=\frac{D}{1-D}\left[-\frac{Q}{1-Q}
 +\frac{P(1+1/Q)}{1-P}\right].
\end{align*}
They follow from $PD^7=Q^2$ and the two rearrangements
\[
 P-Q^2=(1+Q)(1-Q)-(1-P)
 =P(1+Q)(1-Q)-Q^2(1-P).
\]
These coefficients match exactly the logarithm of $S_5(z;s,d)/[S_5(z;t,d)S_5(z+s;t,d)]$. A common region exists: choose $\Im s>6\Im d$ and $0<\Im z<\Im s-6\Im d$. Then $\Im t=2\Im s-7\Im d>0$ and every series is absolutely convergent. Exponentiating proves the identity there. The positive-parameter cone $\Im d>0,\ 2\Im s>7\Im d$ is connected; meromorphic continuation in parameters and argument proves \eqref{eq:p7} throughout. The real shift four contributes no extra exponential.

The raw identity $P_7(z)=F(\beta z)$ is false. At $z=\tau=-t$, $F$ has a zero of order four while the other six shifted factors are nonzero: their $t$ coefficients have nonintegral residues $-k/7$. The denominator argument of positive-period $F(\beta\tau)$ has mixed-sign coordinates $(-1,4)$, since $\beta\tau=2\tau+4s+1$, and is regular and nonzero. Thus $\operatorname{ord}_{\tau}R=4$. An overall scalar cannot repair a nonzero divisor order.
